% Section : sûreté, fiabilité et certification (axe safety-certification + détection de personnes)
\section{Sûreté, fiabilité et certification}
\label{sec:surete}

Les sections précédentes jugent un accélérateur à son débit, à sa puissance et à sa
précision. Pour un \glslink{yolo}{YOLO} qui surveille des personnes dans une zone industrielle dangereuse,
la question change : il ne suffit plus de dire combien d'images par seconde le système
traite, il faut dire quelle preuve on peut fournir qu'il détecte les personnes avec un taux
de défaillance acceptable, et que cette preuve vaut pour le réseau quantifié et compilé qui
tourne sur la carte. \enquote{Sécuritaire} s'entend ici d'abord au sens de la sûreté de
fonctionnement (\emph{safety}) ; la sécurité face aux attaques (\emph{security}) est un
volet secondaire. Le corpus compte 29 fiches sur cet axe. Trois seulement mesurent quelque
chose sur un FPGA, et aucune ne porte sur un détecteur d'objets.

\subsection{Cadres de certification et guides d'assurance}

\paragraph{L'EASA : learning assurance et cycle en W.} Le rapport CoDANN
\citep{2020-easa-concepts-design-assurance-neural-networks} pose le vocabulaire repris
ensuite : un cycle de vie en W qui ajoute au cycle en V la gestion des données, de
l'apprentissage et la vérification du modèle d'inférence \src{§6.1}, la séparation entre
fonction et moniteur (\emph{Doer/Checker}) avec un discriminateur hors distribution
\src{§6.6}, et l'allocation d'objectifs quantitatifs de $10^{-7}$ et $10^{-9}$ par heure de
vol \src{§9.2.5, Fig.~9.3}. Deux limites comptent pour la thèse. Les bornes de
généralisation théoriques sont jugées inutilisables pour de gros réseaux, si bien que les
garanties reposent sur un échantillonnage i.i.d.\ du domaine opérationnel, lui-même
difficile à démontrer \src{§5.3.5, §6.2.8}. Le matériel est hors périmètre : la conformité
DO-254 des GPU, ASIC et FPGA \enquote{will not be addressed} \src{§4.2.3}, et la
vérification de l'étape d'implémentation est laissée à un travail futur \src{§6.5}.

CoDANN~II \citep{2024-easa-concepts-design-assurance-neural-networks} traite justement le
passage du modèle entraîné au modèle déployé. Il distingue équivalence ponctuelle,
proximité ponctuelle et équivalence en moyenne entre les deux \src{§3.6.1}, et propose
d'évaluer la borne de généralisation directement sur le modèle transformé (quantifié,
élagué) dans l'environnement d'inférence \src{§3.7.3}, avec qualification des outils selon
DO-178C §12.2 \src{§3.7.1}. Il exige que le moniteur hors distribution reçoive le même
niveau d'assurance que la fonction principale \src{§5.2.2}. Il argumente enfin pour le FPGA
contre le GPU : logique personnalisée ou Soft IP plus visibles, alors que le GPU repose sur
des pilotes, un microcode et un ordonnancement propriétaires \src{§3.2.2, Annexe~B}. Cet
argument est une position d'entreprise (Daedalean), non étayée par des mesures ; les
sections sur la sûreté du système et les \glslink{dal}{DAL} sont retirées de la version publique.

Le concept paper Issue~02 \citep{2024-easa-concept-paper-guidance-level-1} est le
document d'autorité le plus concret du corpus, mais il reste non contraignant (moyens de
conformité \enquote{anticipated}). Quatre points touchent directement le sujet :
\begin{itemize}
  \item l'apprentissage en ligne n'est pas accepté : le modèle est gelé à l'approbation
    \src{§C.1, MOC LM-02} ;
  \item les FPGA et ASIC relèvent des guides existants (AMC 20-152A), alors que les GPU et
    autres accélérateurs d'IA demanderont des guides propres \src{MOC IMP-05} ;
  \item la quantification qui change le comportement doit être traitée dans l'apprentissage
    \src{LM-06} ; en implémentation, seules les optimisations sans effet sur le
    comportement sont admises \src{MOC IMP-04-2}, et la préservation des propriétés après
    conversion doit être vérifiée, méthodes formelles \enquote{expected to be necessary}
    \src{IMP-06} ;
  \item le système surveille que ses entrées restent dans l'ODD et en distribution, et que
    ses sorties respectent des bornes de performance et de confiance \src{EXP-04 à EXP-07},
    avec repli sur un système classique \src{SRM-02}.
\end{itemize}
Le seul cas d'usage en vision (Visual Landing System) tourne sur un GPU COTS non certifié
\src{§F.2.1.3.3}, et le document ne dit rien de l'équivalence numérique entre modèle
entraîné et bitstream, ni des SEU.

\paragraph{Guides académiques.} AMLAS
\citep{arxiv-2102.01564-guidance-assurance-machine-learning-autonomous}\pp{} propose six
étapes (périmètre, exigences, données, apprentissage, vérification, déploiement), chacune
close par un patron d'argument en GSN. Ses exemples sont proches du cas d'application,
comme l'exigence \enquote{identifier toute personne présente avec une précision d'au moins
0,93} \src{Stage~2, Ex.~7}. Le matériel n'y apparaît qu'à l'intégration, pour vérifier que la
plateforme cible convient au modèle \src{Stage~6, Ex.~45}.
\citet{arxiv-1905.04223-assuring-machine-learning-lifecycle-desiderata}\pp{} structurent le
cycle de vie en quatre étapes et listent les desiderata de chacune ; ils incluent la surveillance des internes du modèle contre les SEU \src{§7.2.2} et le test sur
la cible, parce que la représentation numérique change entre étapes \src{§7.4.1}. Le livre
blanc du projet DEEL \citep{arxiv-2103.10529-white-paper-machine-learning-certified}\pp{}
formule sept défis (évaluation probabiliste, résilience, spécifiabilité, données,
explicabilité, robustesse, vérifiabilité) \src{§4}. Deux idées y sont directement
transposables : une perte asymétrique où la non-détection d'un piéton est catastrophique et
la fausse détection seulement \enquote{major} \src{§4.2.4.2}, et la conversion d'un taux
d'erreur par image en taux de défaillance par heure \src{§4.2.4.4}. L'implémentation est
hors périmètre \src{§1.3}. Les surveys de
\citet{2020-schwalbe-survey-methods-safety-assurance-machine},
\citet{arxiv-1709.02435-analysis-iso-26262-machine-learning}\pp{} et
\citet{arxiv-1812.08342-survey-safety-trustworthiness-deep-neural}\pp{} partent tous
d'ISO~26262 ou de l'automobile. Le deuxième note qu'une fonction de perception comme
\enquote{reconnaître un piéton} est peut-être inspécifiable, et recommande pour les ASIL
élevés des reconnaisseurs redondants entraînés sur des modèles et des données différents
\src{§III-C}.

\paragraph{Un processus appliqué à un YOLO de surveillance.}
\citet{arxiv-2406.12362-certified-ml-object-detection-surveillance} appliquent en partie le
processus ARP6983/ED-324, encore en rédaction, à la détection de drones intrus par
\glslink{yolov3-tiny}{YOLOv3-tiny} sur \glslink{agx-xavier}{Jetson AGX Xavier}. Ils dérivent un ODD du constituant ML (taille des
objets, position, luminosité) \src{§II}, corrigent le biais de position du jeu de données
\src{§III}, découpent l'image 1920$\times$1080 en tuiles \src{Table~I} et réécrivent le
GEMM en CUDA pour éviter cuBLAS, fermé et difficile à analyser \src{§V}. Les exigences
citées sont une détection supérieure à 80\,\% et une latence inférieure à 50\,\glslink{u-ms}{ms}
\src{§I-A} ; la latence obtenue est de 15\,\glslink{u-ms}{ms} en FP16, sans dire si elle porte sur une tuile
ou sur l'image \src{§V-B}. Aucune métrique de détection n'est publiée (confidentialité,
§I-A), et la traçabilité s'arrête aux pilotes NVIDIA \src{§VI}. C'est le travail du corpus
le plus proche du cas d'application, mais sans FPGA.

\paragraph{Normes industrielles : citées, pas traitées.} Aucune fiche ne traite l'IEC 61508,
l'ISO 13849, l'IEC 62061 ou l'ISO/IEC TR~5469 pour elles-mêmes. L'IEC 61508 n'apparaît que
comme contexte \citep[§1, §2.3]{2017-li-understanding-error-propagation-deep-learning},
avec la cible de moins de 10~\glslink{fit}{FIT} tirée d'ISO~26262, reprise pour l'ASIL~D par
\citet[§III-C]{2024-rech-artificial-neural-networks-space-safety}. La revue de
\citet{2018-halme-review-vision-safety-systems-human} s'appuie sur l'ISO 10218-1/2 et
l'ISO/TS 15066 \src{§2} ; sur 15 méthodes de vision pour la collaboration humain-robot,
seul SafetyEYE est commercial et certifié \src{Table~1}. Les correspondances entre les
objectifs EASA ou AMLAS et les niveaux \glslink{sil}{SIL} ou PL ne figurent que dans les gaps rédigés par
les fiches \deduit.

\subsection{Fiabilité matérielle : SEU, radiations, injection de fautes}

\citet{2017-li-understanding-error-propagation-deep-learning} caractérisent par simulation
(pas de RTL, pas de faisceau) la propagation de bit-flips transitoires dans un accélérateur
de DNN. Le taux de corruptions silencieuses dépend fortement du type de données : 7,19\,\%
en virgule fixe 32b\_rb10 contre 0,38\,\% en FLOAT pour \glslink{alexnet}{AlexNet} \src{§5.1.1}, et seuls les
bits de poids fort comptent \src{§5.1.2}. Un détecteur fondé sur des bornes d'activation par
couche (SED) atteint 90,21\,\% de précision et 92,5\,\% de rappel \src{§6.2} (l'introduction
donne d'autres chiffres, §1) et ramène le \glslink{fit}{FIT} d'\glslink{eyeriss}{Eyeriss} de 8,55 à 0,35 en FLOAT
\src{§6.2}. Il est aveugle aux corruptions qui restent dans la plage apprise.

Les deux articles de Libano et al. sont les seuls du corpus à étudier les SEU de la mémoire
de configuration d'un FPGA SRAM.
\citet{2019-libano-selective-hardening-neural-networks-fpgas} isolent chaque couche dans un
Pblock, injectent des fautes dans la mémoire de configuration et valident sous neutrons
(LANSCE, ChipIR) \src{§III}. Sur un ANN de 19 poids (\glslink{xc7z020}{XC7Z020}), le TMR de la seule couche la
plus vulnérable réduit la probabilité d'erreur de 64,2\,\%, contre 85,1\,\% pour le TMR
complet \src{§V-C, Fig.~7}. Sur un CNN MNIST de 1\,848 poids (ZU9EG), il masque environ
40\,\% des fautes pour 8\,\% de ressources en plus ; le TMR complet ne tient pas dans la
puce \src{§V-C}. Les \glslink{bram}{BRAM}, où résident souvent les poids, sont exclues de l'injection
\src{§III-A}. \citet{2020-libano-understanding-impact-quantization-accuracy-radiation}
binarisent les couches convolutives du même type de réseau : la section efficace estimée
baisse de 39\,\%, mais la part d'erreurs critiques passe de 26 à 38\,\% \src{§IV-B}. Leur
indicateur MEBF, qui combine inexactitude du modèle et radiations, montre qu'au niveau de la
mer l'inexactitude domine \src{§IV-B, Fig.~5} ; les auteurs notent eux-mêmes que ce
résultat, obtenu pour une inférence de 233\,\glslink{u-ns}{ns}, ne se transpose pas tel quel à un détecteur
\src{§V}.

\citet{arxiv-1912.07394-efficient-error-tolerant-quantized-neural} injectent des fautes
\enquote{stuck-at} par canal dans des \glslink{qnn}{QNN} FINN sur \glslink{zcu104}{ZCU104}. Sur \glslink{cnv}{CNV} (CIFAR-10), un seul
canal bloqué coûte jusqu'à environ 10,8~points de précision \src{Table~II}. Répliquer
seulement les canaux critiques coûte 49,87\,\% de \glslink{mac}{MAC} en plus pour une chute tolérée de
2\,\%, et 173,47\,\% pour 0,5\,\%, contre 200\,\% pour le TMR complet \src{Table~III}. Les
SEU de configuration sont renvoyés à un travail futur \src{§VII}. FT-ClipAct
\citep{arxiv-1912.00941-ft-clipact-resilience-analysis-deep}\pp{} remplace \glslink{relu}{ReLU} par une
activation écrêtée à un seuil appris : à un taux de fautes de $5\times10^{-7}$ dans les
poids, \glslink{alexnet}{AlexNet} garde 69,36\,\% de précision contre 51,16\,\% sans protection \src{§V-B,
Fig.~7a}, en simulation PyTorch seulement.

Deux surveys cadrent ces travaux. \citet{2023-ruospo-survey-deep-learning-resilience-assessment}
classent l'évaluation en quatre familles (simulation logicielle, simulation matérielle,
plateforme, radiation) \src{Table~1} ; la simulation RTL d'un petit CNN demande environ
25\,min par inférence \src{Discussion}, et les résultats d'injection dans la mémoire de
configuration varient d'un composant à l'autre et avec la température \src{Open
challenges}. \citet{2024-rech-artificial-neural-networks-space-safety} rassemble les
résultats sous faisceau, surtout sur GPU. Il définit une erreur critique en détection par
un seuil de Jaccard sur les boîtes \src{§II-C} et recense des mitigations peu coûteuses :
ABFT jusqu'à 85\,\% des erreurs critiques détectées, cohérence temporelle entre images
environ 70\,\% \src{§VI-A}, RP-DWC 75\,\% pour moins de 20\,\% de surcoût \src{§VI-B}. Ces
chiffres viennent d'autres travaux et ne sont pas comparables entre eux.

Toutes les mesures portent sur des classifieurs, de 19 poids à \glslink{vgg16}{VGG-16}. Hormis le seuil de
Jaccard de Rech, la criticité d'une erreur de détecteur (personne manquée contre fausse
alarme) n'est pas définie, et le passage d'un \glslink{fit}{FIT} mesuré à un objectif \glslink{sil}{SIL} ou PL n'est
traité nulle part \deduit.

\subsection{Vérification formelle}

Reluplex \citep{arxiv-1702.01135-reluplex-efficient-smt-solver-verifying}\pp{} étend le
simplexe aux contraintes \glslink{relu}{ReLU}. Sur les 45 réseaux ACAS~Xu de 300 \glslink{relu}{ReLU}, qui remplacent une
table de plus de 2\,Go par moins de 3\,\glslink{u-mo}{Mo} \src{§5}, la propriété $\varphi_1$ est prouvée sur
41 réseaux en 394\,517\,\glslink{u-s}{s} au total, avec 4 timeouts \src{Table~2}, et $\varphi_8$ révèle un
contre-exemple. Le solveur travaille en flottant et ne garantit pas la correction
\src{§4} : il vérifie le modèle idéalisé, pas le réseau déployé \deduit.
\citet{arxiv-1709.06662-verifying-properties-binarized-deep-neural}\pp{} encodent
exactement un \glslink{bnn}{BNN} en SAT, ce qui permet aussi de vérifier l'équivalence de deux réseaux
\src{§5}. Sur un réseau dense de cinq blocs, avec un timeout de 300\,\glslink{u-s}{s}, la robustesse à
$\epsilon = 3$ n'est certifiée que pour au plus 21 des 200 images de MNIST-back-image
\src{Table~2}. \citet{arxiv-2012.08185-scalable-verification-quantized-neural-networks}
vérifient bit à bit des réseaux en virgule fixe ; ils montrent que le problème est
PSPACE-difficile avec une spécification bit-vector \src{Théorème~1}, et annoncent un gain
allant jusqu'à trois ordres de grandeur sur les outils existants \src{Abstract}, sur un
réseau dense de 52\,650 paramètres. Leurs tableaux de résultats sont mal extraits ; la
fiche ne garantit pas les chiffres cellule par cellule.

L'écart d'échelle est de plusieurs ordres de grandeur : le survey de
\citet[§4.2, §4.3]{arxiv-1812.08342-survey-safety-trustworthiness-deep-neural}\pp{} situe
les méthodes à bornes approchées à quelques milliers, au plus 10\,000, neurones cachés, et
le concept paper EASA juge les méthodes formelles limitées à des évaluations locales
\citep[§C.3.1.6]{2024-easa-concept-paper-guidance-level-1}. Prouver une propriété d'un \glslink{yolo}{YOLO}
complet n'est pas à portée. Deux usages restent plausibles \deduit : vérifier des
sous-composants (une tête, un moniteur binaire), et vérifier l'équivalence entre le réseau
quantifié de référence et ce que calcule le circuit, ce qu'aucun article du corpus ne fait.

\subsection{Surveillance à l'exécution et incertitude}

Les moniteurs du corpus sont presque tous évalués sur des classifieurs :
\begin{itemize}
  \item la probabilité softmax maximale (MSP)
    \citep{arxiv-1610.02136-baseline-detecting-misclassified-out-distribution} sépare
    CIFAR-10 de SUN avec une \glslink{auroc}{AUROC} de 95 \src{Table~2}, mais tombe à 61 sur une tâche de NLP
    presque dans la distribution \src{Table~8} ;
  \item les patrons d'activation on/off stockés dans un BDD
    \citep{arxiv-1809.06573-runtime-monitoring-neuron-activation-patterns}\pp{} : sur
    GTSRB, à $\gamma = 3$, le moniteur se tait 95,42\,\% du temps et, quand il alerte,
    l'image est mal classée dans 54,54\,\% des cas \src{Table~I} ; les BDD limitent
    pratiquement la couche surveillée à environ 200 neurones \src{§II} ;
  \item les boîtes d'abstraction sur les couches cachées
    \citep{arxiv-1911.09032-outside-box-abstraction-monitoring-neural} : une entrée
    d'entraînement n'est jamais rejetée \src{Prop.~1}, la mémoire ne dépend pas de la taille
    du jeu \src{§4.3} et le test prend moins d'une milliseconde par entrée sur un portable
    \src{§5} ;
  \item SafeML \citep{arxiv-2005.13166-safeml-safety-monitoring-machine-learning}\pp{}
    compare des distributions empiriques sur un tampon d'échantillons ; sur un cas jouet,
    il estime une exactitude minimale de 0,772 là où la vraie vaut 0,508 \src{Table~1},
    c'est-à-dire qu'il se trompe dans le sens dangereux.
\end{itemize}

Trois travaux portent sur la détection d'objets. \glslink{gaussian-yolov3}{Gaussian YOLOv3}
\citep{arxiv-1904.04620-gaussian-yolov3-accurate-fast-object}\pp{} prédit une variance par
coordonnée de boîte pour 0,04\,\% de \glslink{flops}{FLOPs} en plus \src{§3.1} et réduit les faux positifs
de 41,40\,\% sur KITTI au seuil 0,5 \src{Table~4} ; c'est une incertitude aléatoire de
localisation, sans calibration ni OOD, mesurée sur GTX~1080~Ti.
\citet{arxiv-2011.10671-review-comparative-probabilistic-object-detection} comparent les
méthodes d'incertitude sur RetinaNet : la \glslink{map}{mAP} ne mesure pas la qualité de l'incertitude, la
modélisation directe de la variance est essentielle, et MC-Dropout ou les ensembles
n'apportent pas assez pour leur coût (2 à 3~\glslink{fps}{FPS} pour 10 passes de MC-Dropout) \src{§IV-D,
§V-A}. Sous décalage de domaine (Lyft), la \glslink{map}{mAP} s'effondre vers 5 à 6 pour toutes les
variantes \src{Table~IV, corps possiblement décalé à l'extraction}. Enfin,
\citet{arxiv-2304.06052-confident-object-detection-conformal-prediction}\pp{} appliquent la
prédiction conforme à \glslink{yolov5m}{YOLOv5m}, \glslink{detr}{DETR} et DiffusionDet sur 3\,414 images de signaux
ferroviaires : la couverture empirique des boîtes va de 0,92 à 1,00 pour une cible de 0,9
\src{Table~4}. La garantie ne porte que sur les vrais positifs \src{§4.1} : une personne
manquée n'est pas couverte, et l'hypothèse d'échangeabilité ne tient pas sur une longue
durée \src{§6}.

Aucun de ces moniteurs n'est implanté sur FPGA, ni évalué sur des personnes en milieu
industriel. Plusieurs se réduisent à des comparaisons ou à des additions (seuil sur la MSP,
test d'appartenance à une boîte, marge conforme), ce qui les rend a priori peu coûteux en
logique \deduit ; leur coût réel et leur effet sur un réseau INT8 restent à mesurer.

\subsection{Logiciel prédictible et traçabilité de l'implémentation}

ACETONE \citep{2022-de-albuquerque-silva-acetone-predictable-programming-framework-ml}
génère du C séquentiel, à allocation statique, depuis un réseau entraîné, pour la
traçabilité et le calcul de \glslink{wcet}{WCET} exigés par DO-178C \src{§1, §3}. L'écart numérique au
modèle est de l'ordre de $10^{-6}$ en FP32 \src{§5}, mais la préservation de la sémantique
est testée, pas prouvée. Le \glslink{wcet}{WCET} est calculé par OTAWA sur un lpc2138, non représentatif du
\glslink{cortex-a15}{Cortex-A15} mesuré \src{§4.2}, et l'analyse échoue faute de mémoire pour \glslink{alexnet}{AlexNet} et la
variante déroulée \src{§5}. Rien d'équivalent n'existe dans le corpus pour un flot HLS ou
RTL \deduit.

Côté FPGA, les accélérateurs \glslink{yolo}{YOLO} du corpus ne traitent pas la sûreté. Sur \glslink{kv260}{KV260},
\citet{2026-fata-low-cost-fpga-enhanced-cnn} renvoient la tolérance aux fautes à un travail
futur \src{§VI}, et leur latence varie selon la vidéo \src{Fig.~14}. Sur \glslink{zcu104}{ZCU104},
\citet{arxiv-2607.13106-no-attention-no-problem-dpu}\pp{} détectent des personnes sur un jeu
CCTV non publié ; la \glslink{map5095}{mAP@0.5:0.95} passe de 0,75 sur CPU à 0,71 sur le FPGA en INT8 pour
\glslink{yolov11n}{YOLOv11n}, et de 0,64 à 0,34 pour \glslink{yolov8n}{YOLOv8n} sur KITTI \src{Table~II}. Un tel écart, imprévisible
d'un modèle à l'autre, est exactement ce que l'objectif IMP-06 demanderait de borner
\deduit. Le constat rejoint celui de la section~\ref{sec:critique} : la précision n'est
presque jamais remesurée sur la cible.

\subsection{Détection de personnes et robustesse}

La revue de \citet{2018-halme-review-vision-safety-systems-human} donne le cadre fonctionnel
du cas d'application (arrêt de sécurité, surveillance de vitesse et de distance de
séparation) et ses difficultés : occultations, sensibilité des caméras 2D à la lumière et à
la poussière, latences de perception non traitées \src{§4}. Elle n'évalue aucun détecteur de
personnes appris. CrowdHuman \citep{arxiv-1805.00123-crowdhuman-benchmark-detecting-human-crowd}\pp{}
mesure la difficulté de l'occultation : 22,64 personnes par image en moyenne \src{Table~2},
et un FPN atteint un taux d'échec moyen (mMR) de 50,42 sur le corps entier contre 10,08 sur
Caltech \src{Table~7}. Ce sont des images du web, sans scène industrielle ni EPI. Le survey
de ré-identification \citep{arxiv-2001.04193-deep-learning-person-re-identification}\pp{}
note une forte chute des performances entre jeux de données et sous attaque adverse
\src{§2.4.2}, sans aucune mesure embarquée. Pour la sécurité,
\citet{arxiv-1904.08653-fooling-automated-surveillance-cameras-adversarial}\pp{} montrent
qu'un patch imprimé fait tomber le rappel de \glslink{yolov2}{YOLOv2} sur Inria de 100 à 26,46\,\%
\src{Table~1} ; le 100\,\% vient de ce que les vérités terrain sont produites par le même
détecteur \src{§4}. L'effet de la quantification et de la compilation FPGA sur cette
vulnérabilité n'est pas étudié.

\subsection{Bilan}

\begin{table}[!htbp]
\centering\footnotesize
\caption{Familles de garanties dans le corpus : cible matérielle, modèle étudié et limite
principale pour un YOLO certifiable sur FPGA.}
\label{tab:surete}
\begin{tabularx}{\textwidth}{@{}P{2.5cm}P{3cm}P{2.2cm}P{2.3cm}L@{}}
\toprule
Famille & Travaux & Cible & Modèle & Limite principale \\
\midrule
Cadres et guides & EASA\,\textsuperscript{a}, AMLAS, DEEL, Ashmore & aucune & générique & matériel hors périmètre ou renvoyé à AMC 20-152A ; aviation ou automobile \\
Processus appliqué & ARP6983\,\textsuperscript{b} & \glslink{agx-xavier}{Jetson AGX Xavier} & \glslink{yolov3-tiny}{YOLOv3-tiny} & pas de FPGA ; aucune métrique de détection publiée \\
Injection de fautes, radiations & Li, Libano\,\textsuperscript{c}, Gambardella & simulateur ; \glslink{zynq}{Zynq}, ZU9EG, \glslink{zcu104}{ZCU104} & classifieurs (19 poids à NiN) & aucun détecteur ; pas de lien FIT--SIL \\
Mitigation logicielle & FT-ClipAct, Rech\,\textsuperscript{d} & simulation ; GPU surtout & classifieurs & transfert vers FPGA quantifié non fait \\
Vérification formelle & Reluplex, \glslink{bnn}{BNN}\,\textsuperscript{e}, \glslink{qnn}{QNN}\,\textsuperscript{f} & CPU (solveurs) & denses, $\leq$ 52\,650 paramètres & échelle ; pas d'équivalence avec le circuit \\
Moniteurs OOD & MSP, patrons, boîtes, SafeML & CPU, GPU & classifieurs & jamais sur détecteur ni sur FPGA \\
Incertitude, conforme & \glslink{gaussian-yolov3}{Gaussian YOLOv3}, POD\,\textsuperscript{g}, conforme\,\textsuperscript{h} & GPU de bureau & \glslink{yolov3}{YOLOv3}, RetinaNet, \glslink{yolov5m}{YOLOv5m} & faux négatifs non couverts ; coût FPGA inconnu \\
Code prédictible & ACETONE & CPU ARM & MLP, CNN FP32 & pas de HLS/RTL ni de quantification \\
Personnes, robustesse & Halme, CrowdHuman, Re-ID, patch\,\textsuperscript{i} & aucune ou GPU & FPN, \glslink{yolov2}{YOLOv2} & ni scène industrielle ni cible embarquée \\
\bottomrule
\end{tabularx}

\smallskip\raggedright\footnotesize
\textsuperscript{a}\,\citet{2020-easa-concepts-design-assurance-neural-networks},
\citet{2024-easa-concepts-design-assurance-neural-networks},
\citet{2024-easa-concept-paper-guidance-level-1} ;
\textsuperscript{b}\,\citet{arxiv-2406.12362-certified-ml-object-detection-surveillance} ;
\textsuperscript{c}\,\citet{2019-libano-selective-hardening-neural-networks-fpgas},
\citet{2020-libano-understanding-impact-quantization-accuracy-radiation} ;
\textsuperscript{d}\,\citet{2024-rech-artificial-neural-networks-space-safety} ;
\textsuperscript{e}\,\citet{arxiv-1709.06662-verifying-properties-binarized-deep-neural}\pp ;
\textsuperscript{f}\,\citet{arxiv-2012.08185-scalable-verification-quantized-neural-networks} ;
\textsuperscript{g}\,\citet{arxiv-2011.10671-review-comparative-probabilistic-object-detection} ;
\textsuperscript{h}\,\citet{arxiv-2304.06052-confident-object-detection-conformal-prediction}\pp ;
\textsuperscript{i}\,\citet{arxiv-1904.08653-fooling-automated-surveillance-cameras-adversarial}\pp.
\end{table}

Le tableau~\ref{tab:surete} montre des briques qui ne se rejoignent pas. Les cadres
disent quoi prouver mais pas comment le faire sur un FPGA ; les méthodes de fiabilité,
de vérification et de surveillance savent faire sur un petit classifieur, presque jamais
sur un détecteur, et jamais sur un détecteur implanté sur FPGA. Pour un \glslink{yolo}{YOLO} certifiable
sur FPGA, il manquerait au minimum \deduit :
\begin{itemize}
  \item une correspondance entre les objectifs EASA ou AMLAS et les normes industrielles
    (IEC 61508, ISO 13849, ISO/IEC TR~5469) ;
  \item une métrique de criticité propre à la détection de personnes, reliée à un taux de
    défaillance par heure, puis à un \glslink{sil}{SIL} ou un PL ;
  \item une campagne d'injection de fautes et un durcissement sélectif sur un \glslink{yolo}{YOLO} quantifié,
    \glslink{bram}{BRAM} comprises ;
  \item une preuve, ou au moins une borne mesurée, de l'écart entre le modèle de référence et
    le circuit déployé ;
  \item un moniteur d'exécution implanté sur la même puce, au coût et à la latence mesurés ;
  \item un jeu de données de personnes en zone industrielle, sur lequel évaluer le tout.
\end{itemize}
L'EASA ajoute une tension avec l'axe apprentissage continu de ce document : le modèle doit
être gelé à l'approbation \citep[MOC LM-02]{2024-easa-concept-paper-guidance-level-1}. Ces
manques sont repris parmi les gaps de la section~\ref{sec:gaps}.
